% !TeX root = ../../ONMIPA Alin.tex
% Sumber utama: Zhang (2009), Bab 3; Zhang (2011), bagian 3.1--3.4.
\chapter{Dekomposisi Matriks}
\label{ch:dekomposisi}

Bab sebelumnya menyediakan nilai eigen, ruang eigen, dan polinom minimal.
Sekarang kita memilih basis yang menyederhanakan matriks: mula-mula
basis vektor eigen, lalu basis rantai Jordan, dan akhirnya basis
ortonormal pada dekomposisi Schur serta spektral.

\section{Diagonalisasi}
\label{eigen:diagonalisasi}

\begin{definisi}[][Matriks yang dapat didiagonalkan][dekomp:def-diagonal]
    Matriks $A\in\Mat_n(\F)$ dapat didiagonalkan di atas $\F$ jika
    $A=PDP^{-1}$ dengan $D$ diagonal dan $P$ invertibel, keduanya
    berentri dalam $\F$.
\end{definisi}

\begin{teorema}[][Kriteria diagonalisasi][eigen:kriteria-diagonal]
    Untuk $A\in\Mat_n(\F)$, pernyataan berikut ekuivalen:
    \begin{enumerate}[label=(\roman*),nosep]
        \item $A$ dapat didiagonalkan di atas $\F$;
        \item terdapat basis $\F^n$ yang seluruh anggotanya vektor eigen $A$;
        \item $\chi_A$ terurai menjadi faktor linear di atas $\F$ dan
            $g_\lambda=a_\lambda$ untuk setiap nilai eigen $\lambda$.
    \end{enumerate}
\end{teorema}
\begin{proof}
    Jika $P=[v_1\ \cdots\ v_n]$ dan
    $D=\diag(\lambda_1,\ldots,\lambda_n)$, persamaan $AP=PD$
    ekuivalen dengan $Av_i=\lambda_iv_i$ untuk setiap $i$.
    Ini membuktikan (i)--(ii).
    Ruang-ruang eigen berbeda berjumlah langsung dan dimensinya
    tidak melebihi multiplisitas aljabarnya.
    Karena $\sum_\lambda a_\lambda=n$ ketika polinomnya terurai,
    basis vektor eigen ada tepat ketika setiap batas tersebut tercapai.
\end{proof}

Adanya $n$ nilai eigen berbeda cukup untuk diagonalisasi, tetapi
tidak perlu: matriks identitas mempunyai satu nilai eigen dan tetap
diagonal. Sebaliknya, akar berulang mengharuskan pemeriksaan dimensi
ruang eigen; akar berulang sendiri tidak memutuskan jawabannya.

\begin{exam}[Parameter pada ruang eigen berulang]
    Tentukan kapan
    \[
        A_t=\begin{pmatrix}1&t&0\\0&1&0\\0&0&2\end{pmatrix}
    \]
    dapat didiagonalkan.
\end{exam}
\begin{solution}
    Untuk setiap $t$, $\chi_{A_t}(x)=(x-1)^2(x-2)$.
    Persamaan $(A_t-I)v=0$ memberi $tv_2=0$ dan $v_3=0$.
    Jadi $g_1=2$ jika $t=0$, tetapi $g_1=1$ jika $t\ne0$.
    Karena $g_2=1$, matriks dapat didiagonalkan tepat ketika $t=0$.
    Contoh ini juga menunjukkan bahwa polinom karakteristik yang sama
    dapat dimiliki oleh matriks yang tidak serupa.
\end{solution}

\begin{exam}[Pola diagonal dan luar diagonal konstan]
    Misalkan $A$ berukuran $n\ge2$, dengan entri diagonal $a$ dan
    semua entri luar diagonal $b$. Tentukan ruang eigen dan pangkat $A$.
\end{exam}
\begin{solution}
    Dengan $e=(1,\ldots,1)^T$ dan $J_n=ee^T$,
    \[
        A=(a-b)I+bJ_n,\qquad J_ne=ne.
    \]
    Pada ruang $H=\{x:e^Tx=0\}$, berlaku $J_nx=0$.
    Karena $\F^n=\operatorname{span}\{e\}\oplus H$, matriks $A$
    bekerja sebagai perkalian dengan $a+(n-1)b$ pada ruang pertama
    dan $a-b$ pada ruang kedua. Jika $b\ne0$, kedua nilai ini berbeda
    dan multiplisitasnya $1$ serta $n-1$. Jika $b=0$, $A=aI$ dan
    ruang eigen untuk $a$ adalah seluruh $\F^n$.

    Definisikan $P=J_n/n$ dan $Q=I-P$.
    Karena $P^2=P$, $Q^2=Q$, dan $PQ=QP=0$,
    \[
        A^k=\bigl(a+(n-1)b\bigr)^kP+(a-b)^kQ\qquad(k\ge1).
    \]
    Untuk $k=0$, $A^0=P+Q=I$.
    Jika kedua skalar tak nol, invers diperoleh dengan mengganti
    pangkat $k$ oleh $-1$.
\end{solution}


\begin{exam}[Menyusun matriks perubahan basis]
    Diagonalkan
    \[
        A=\begin{pmatrix}1&2&2\\2&1&2\\2&2&1\end{pmatrix}.
    \]
\end{exam}
\begin{solution}
    Karena $A=-I+2J_3$, nilai eigennya $5,-1,-1$.
    Pilih $v_1=(1,1,1)^T$, $v_2=(1,-1,0)^T$, dan $v_3=(1,1,-2)^T$.
    Ketiganya bebas dan memenuhi
    \[
        P=\begin{pmatrix}1&1&1\\1&-1&1\\1&0&-2\end{pmatrix},\qquad
        AP=P\diag(5,-1,-1).
    \]
    Jadi $P^{-1}AP=\diag(5,-1,-1)$.
    Jika tiap kolom dinormalisasi, diperoleh matriks ortogonal dan
    diagonalisasi yang sekaligus mempertahankan hasil kali dalam.
\end{solution}

Medan skalar tetap penting. Matriks real
$R=\begin{pmatrix}0&-1\\1&0\end{pmatrix}$ mempunyai polinom
$t^2+1$, sehingga tidak mempunyai vektor eigen real.
Di atas $\C$, nilai eigennya $i,-i$ berbeda dan $R$ dapat didiagonalkan.

\section{Polinom Minimal dan Proyektor}
\label{dekomp:polinom}

Polinom minimal dari \autoref{eigen:minimal} memberikan kriteria
untuk menentukan apakah suatu basis vektor eigen lengkap tersedia.

\begin{teorema}[][Kriteria polinom untuk diagonalisasi][eigen:squarefree]
    Matriks $A$ dapat didiagonalkan di atas $\F$ jika dan hanya jika
    $m_A$ merupakan hasil kali faktor-faktor linear yang berbeda di atas $\F$.
    Lebih umum, cukup ada polinom penganulir $q(A)=0$ yang merupakan
    hasil kali faktor linear berbeda.
\end{teorema}
\begin{proof}
    Jika $A$ diagonal dalam suatu basis, hasil kali
    $\prod_\lambda(t-\lambda)$ atas nilai eigen berbeda menganulir $A$.
    Polinom minimal membagi hasil kali itu dan memuat semua akarnya.

    Sebaliknya, misalkan $q(t)=\prod_{i=1}^s(t-\lambda_i)$ dengan
    $\lambda_i$ berbeda dan $q(A)=0$.
    Bentuk polinom Lagrange
    \[
        \ell_i(t)=\prod_{j\ne i}\frac{t-\lambda_j}{\lambda_i-\lambda_j}.
    \]
    Karena $\sum_i\ell_i(t)=1$ dan $(t-\lambda_i)\ell_i(t)$
    adalah kelipatan $q(t)$, setiap $x\in\F^n$ dapat ditulis
    \[
        x=\sum_i\ell_i(A)x,\qquad
        (A-\lambda_iI)\ell_i(A)x=0.
    \]
    Jadi ruang-ruang eigen merentang seluruh ruang dan berjumlah langsung.
    Gabungan basisnya merupakan basis vektor eigen.
\end{proof}

Akibatnya, $A^2=A$ dan $A^3=A$ menjamin diagonalisasi di atas $\R$
atau $\C$, sebab polinom $t(t-1)$ dan $t(t-1)(t+1)$ mempunyai akar
berbeda. Jika $A^k=I$, $A$ dapat didiagonalkan di atas $\C$, karena
$t^k-1$ tidak mempunyai akar berulang.
Sebaliknya, $A^2=0$ belum menjamin diagonalisasi; matriks nilpoten
yang dapat didiagonalkan harus merupakan matriks nol.

\subsection{Proyektor dan reduksi pangkat}

Jika nilai eigen berbeda dari matriks yang dapat didiagonalkan adalah
$\lambda_1,\ldots,\lambda_s$, definisikan $P_i=\ell_i(A)$.
Dengan memeriksa aksinya pada setiap ruang eigen, diperoleh
\[
    P_i^2=P_i,\qquad P_iP_j=0\ (i\ne j),\qquad
    \sum_iP_i=I,\qquad A=\sum_i\lambda_iP_i.
\]
Maka $f(A)=\sum_i f(\lambda_i)P_i$ untuk setiap polinom $f$.
Proyektor ini belum tentu ortogonal; ortogonalitas akan muncul pada
dekomposisi spektral matriks normal.

\begin{exam}[Identitas kuadratik menentukan semua pangkat]
    Misalkan $A^2-3A+2I=0$. Tentukan $A^k$ dan $A^{-1}$.
\end{exam}
\begin{solution}
    Polinom $(t-1)(t-2)$ mempunyai dua akar berbeda.
    Proyektor yang bersesuaian adalah $P_1=2I-A$ dan $P_2=A-I$;
    salah satunya boleh nol jika hanya satu nilai eigen yang muncul.
    Untuk setiap $k\ge0$,
    \[
        A^k=P_1+2^kP_2=(2-2^k)I+(2^k-1)A.
    \]
    Dari $A(3I-A)=2I$, diperoleh $A^{-1}=(3I-A)/2$.
    Tidak diperlukan perhitungan entri atau vektor eigen satu per satu.
\end{solution}

\section{Diagonalisasi Simultan}
\label{eigen:komutatif}

\begin{prop}[Ruang eigen yang invarian]
    Jika $AB=BA$, maka setiap ruang eigen $E_\lambda(A)$ invarian
    terhadap $B$: $B(E_\lambda(A))\subseteq E_\lambda(A)$.
\end{prop}
\begin{proof}
    Untuk $Ax=\lambda x$, berlaku $A(Bx)=B(Ax)=\lambda Bx$.
    Vektor $Bx$ boleh nol, sehingga pernyataan yang tepat memakai
    ruang eigen, bukan selalu vektor eigen tak nol.
\end{proof}

\begin{teorema}[][Diagonalisasi simultan][eigen:simultan]
    Dua matriks yang masing-masing dapat didiagonalkan di atas $\F$
    dapat didiagonalkan oleh satu perubahan basis yang sama jika dan
    hanya jika keduanya komutatif.
\end{teorema}
\begin{proof}
    Jika keduanya diagonal dalam satu basis, jelas keduanya komutatif.
    Sebaliknya, misalkan $AB=BA$.
    Dekomposisikan $\F^n$ menjadi jumlah langsung ruang eigen $A$.
    Pada setiap ruang ini, pembatasan $B$ tetap dianulir oleh polinom
    minimal $B$, yang mempunyai faktor linear berbeda.
    Maka pembatasan itu dapat didiagonalkan.
    Pilih basis vektor eigen $B$ di setiap ruang eigen $A$, lalu gabungkan;
    setiap anggotanya merupakan vektor eigen untuk kedua matriks.
\end{proof}

\begin{teorema}[][Komutativitas dan spektrum sederhana][eigen:centralizer]
    Jika $A$ mempunyai $n$ nilai eigen berbeda dan $AB=BA$, maka
    terdapat tepat satu polinom $p$ berderajat kurang dari $n$ sehingga
    $B=p(A)$. Khususnya, $B$ juga dapat didiagonalkan.
\end{teorema}
\begin{proof}
    Setiap ruang eigen $A$ berdimensi satu, sehingga $Bv_i=\mu_iv_i$
    pada basis vektor eigen $A$.
    Interpolasi Lagrange memberikan satu polinom berderajat kurang
    dari $n$ dengan $p(\lambda_i)=\mu_i$.
    Maka $p(A)$ dan $B$ mempunyai aksi sama pada suatu basis.
    Keunikannya mengikuti fakta bahwa polinom berderajat kurang dari
    $n$ tidak dapat mempunyai $n$ akar berbeda kecuali polinom nol.
\end{proof}

\section{Bentuk Jordan dan Rantai Vektor Eigen}
\label{eigen:jordan}

Ketika jumlah dimensi ruang eigen kurang dari $n$, vektor eigen biasa
belum membentuk basis. Penggantinya adalah rantai vektor yang
berturut-turut dipetakan ke anggota rantai sebelumnya.

\begin{definisi}[][Blok dan rantai Jordan][eigen:def-jordan]
    Blok Jordan berukuran $s$ untuk $\lambda$ adalah
    \[
        J_s(\lambda)=\lambda I_s+N_s
        =\begin{pmatrix}
            \lambda&1&&0\\&\lambda&\ddots&\\&&\ddots&1\\0&&&\lambda
        \end{pmatrix},
    \]
    dengan $N_s$ mempunyai entri $1$ pada superdiagonal dan nol
    di tempat lain. Untuk $s=1$, bloknya hanya $(\lambda)$.
    Rantai Jordan $v_1,\ldots,v_s$ memenuhi
    \[
        (A-\lambda I)v_1=0,\qquad
        (A-\lambda I)v_j=v_{j-1}\quad(2\le j\le s),\qquad v_1\ne0.
    \]
\end{definisi}

Semua anggota rantai bebas linear. Jika $\sum_jc_jv_j=0$, terapkan
$(A-\lambda I)^{s-1}$ untuk memperoleh $c_sv_1=0$;
lanjutkan dengan pangkat yang lebih rendah untuk mendapatkan semua
koefisien nol. Vektor $v_j$ disebut vektor eigen tergeneralisasi karena
$(A-\lambda I)^jv_j=0$.

\begin{teorema}[][Dekomposisi Jordan][eigen:teorema-jordan]
    Setiap $A\in\Mat_n(\C)$ dapat ditulis
    \[
        A=PJP^{-1},\qquad
        J=\diag\bigl(J_{s_1}(\lambda_1),\ldots,J_{s_r}(\lambda_r)\bigr),
    \]
    dengan $P$ invertibel. Nilai $\lambda_i$ boleh berulang.
    Ukuran blok untuk setiap nilai eigen ditentukan secara unik,
    hingga urutan blok.
\end{teorema}

Secara struktural, dekomposisi ini mula-mula memisahkan ruang menjadi
ruang-ruang eigen tergeneralisasi $\ker(A-\lambda I)^n$.
Pemisahan tersebut mengikuti identitas Bezout untuk faktor-faktor
polinom karakteristik yang saling prima.
Pada ruang untuk $\lambda$, pemetaan $A-\lambda I$ nilpoten dan
dapat disusun dalam basis rantai. Setiap rantai menghasilkan satu
blok Jordan. Pembuktian lengkap keberadaan basis rantai
tidak dirinci di sini.
Untuk perhitungan, ukuran rantainya dapat dibaca dari nulitas,
seperti dibuktikan berikut.

\begin{teorema}[][Membaca blok Jordan dari nulitas][eigen:jordan-nulitas]
    Tetapkan satu nilai eigen $\lambda$ dan tuliskan
    $d_k=\nul(A-\lambda I)^k$, dengan $d_0=0$.
    Jika ukuran blok Jordan untuk $\lambda$ adalah $s_1,\ldots,s_r$, maka
    \[
        d_k=\sum_{j=1}^r\min\{k,s_j\}.
    \]
    Akibatnya,
    \begin{align*}
        \#\{\text{blok berukuran setidaknya }k\}&=d_k-d_{k-1},\\
        \#\{\text{blok berukuran tepat }k\}&=2d_k-d_{k-1}-d_{k+1}.
    \end{align*}
\end{teorema}
\begin{proof}
    Pada satu blok untuk $\lambda$, matriks $A-\lambda I$
    berubah menjadi $N_s$ dan $\nul N_s^k=\min\{k,s\}$.
    Pada blok untuk nilai eigen $\mu\ne\lambda$, matriks
    $J_s(\mu)-\lambda I$ invertibel, sehingga tidak menyumbang kernel.
    Nulitas matriks diagonal blok adalah jumlah nulitas blok-bloknya.
    Satu blok menyumbang $1$ pada $d_k-d_{k-1}$ tepat ketika $s\ge k$;
    selisih dua hitungan berurutan memberi jumlah blok berukuran tepat $k$.
\end{proof}

Dengan demikian, $g_\lambda=d_1$ adalah banyaknya blok untuk $\lambda$,
$a_\lambda=\lim_k d_k$ adalah jumlah ukurannya, dan ukuran blok
terbesar adalah indeks pertama $k$ dengan $d_k=d_{k+1}$.
Ini memberi makna struktural pada stabilisasi kernel yang dibahas
dalam \autoref{rank:deret-pangkat}.

\begin{prop}[Polinom minimal dan blok terbesar]
    \label{eigen:minimal-jordan}
    Jika ukuran blok terbesar untuk nilai eigen $\lambda$ adalah
    $h_\lambda$, maka
    \[
        m_A(t)=\prod_\lambda(t-\lambda)^{h_\lambda}.
    \]
\end{prop}
\begin{proof}
    Satu blok $J_s(\lambda)$ dianulir oleh $(t-\lambda)^s$
    dan tidak oleh pangkat yang lebih kecil karena $N_s^{s-1}\ne0$.
    Faktor polinom yang tidak nol di $\lambda$ memberikan matriks
    segitiga dengan diagonal tak nol, sehingga invertibel pada blok itu.
    Suatu polinom menganulir seluruh matriks tepat ketika menganulir
    setiap blok, sehingga pangkat yang diperlukan untuk tiap faktor
    adalah ukuran blok terbesarnya.
\end{proof}

\begin{exam}[Membangun perubahan basis dari rantai]
    Tentukan bentuk Jordan dan matriks perubahan basis untuk
    \[
        A=\begin{pmatrix}3&-1&0\\1&1&0\\-1&1&3\end{pmatrix}.
    \]
\end{exam}
\begin{solution}
    Polinom karakteristiknya $(t-2)^2(t-3)$.
    Ambil $v_1=(1,1,0)^T$, $v_2=(1,0,1)^T$, dan $w=(0,0,1)^T$.
    Perkalian langsung memberikan
    \[
        (A-2I)v_1=0,\qquad (A-2I)v_2=v_1,\qquad Aw=3w.
    \]
    Kolom-kolom tersebut bebas, sehingga
    \[
        P=\begin{pmatrix}1&1&0\\1&0&0\\0&1&1\end{pmatrix},\qquad
        P^{-1}AP=\begin{pmatrix}2&1&0\\0&2&0\\0&0&3\end{pmatrix}.
    \]
    Urutan $v_1,v_2$ membuat entri $1$ berada di superdiagonal.
    Matriks ini tidak dapat didiagonalkan karena ruang eigen untuk $2$
    hanya berdimensi satu.
\end{solution}

\begin{exam}[Ukuran blok dari beberapa nulitas]
    Diketahui $A\in\Mat_8(\C)$ memenuhi
    \[
        \chi_A(t)=(t-2)^6(t+1)^2,
    \]
    serta
    \[
        \nul(A-2I)=3,\quad\nul(A-2I)^2=5,\quad
        \nul(A-2I)^3=6,\quad\nul(A+I)=1.
    \]
    Tentukan bentuk Jordan dan polinom minimalnya.
\end{exam}
\begin{solution}
    Untuk $\lambda=2$, barisan nulitas dimulai dengan $0,3,5,6,6$.
    Selisihnya $3,2,1,0$, sehingga terdapat tiga blok, dua berukuran
    setidaknya dua, dan satu berukuran setidaknya tiga.
    Ukurannya tepat $3,2,1$.
    Untuk $\lambda=-1$, multiplisitas aljabarnya dua dan jumlah bloknya
    satu, sehingga bloknya berukuran dua. Maka
    \[
        J=\diag\bigl(J_3(2),J_2(2),(2),J_2(-1)\bigr),\qquad
        m_A(t)=(t-2)^3(t+1)^2.
    \]
\end{solution}

\subsection{Pangkat dan polinom pada satu blok}

Karena $\lambda I$ dan $N_s$ komutatif serta $N_s^s=0$,
\[
    J_s(\lambda)^k
    =\sum_{j=0}^{\min\{k,s-1\}}\binom{k}{j}\lambda^{k-j}N_s^j
    \qquad(k\ge0),
\]
dengan faktor berpangkat nol diartikan $1$.
Untuk polinom $f$, ekspansi Taylor yang berhingga memberi
\[
    f(J_s(\lambda))
    =\sum_{j=0}^{s-1}\frac{f^{(j)}(\lambda)}{j!}N_s^j.
\]
Jadi akar berulang pada polinom minimal membuat turunan $f$ ikut
menentukan $f(A)$, selain nilai $f$ pada nilai eigen.

\begin{prop}[Nulitas polinom matriks]
    \label{eigen:nulitas-polinom}
    Misalkan $f$ polinom tak nol dan $q_\lambda$ adalah orde akarnya
    di $\lambda$, dengan $q_\lambda=0$ jika $f(\lambda)\ne0$.
    Setiap blok $J_s(\lambda)$ menyumbang
    $\min\{q_\lambda,s\}$ pada nulitas $f(A)$.
\end{prop}
\begin{proof}
    Tuliskan $f(t)=(t-\lambda)^{q_\lambda}g(t)$ dengan
    $g(\lambda)\ne0$. Pada blok tersebut,
    $f(J_s(\lambda))=N_s^{q_\lambda}g(J_s(\lambda))$.
    Faktor kedua invertibel karena segitiga dengan diagonal
    $g(\lambda)$, sehingga nulitasnya sama dengan nulitas
    $N_s^{q_\lambda}$.
\end{proof}

Pada contoh ukuran delapan di atas, ambil
$f(t)=(t-2)^2(t+1)$. Blok-blok untuk $2$ menyumbang nulitas
$2+2+1=5$, dan blok untuk $-1$ menyumbang $1$.
Jadi $\nul f(A)=6$ dan $\rank f(A)=2$ tanpa perlu mengetahui entri $A$.

\section{Dekomposisi Schur dan Spektral}
\label{eigen:spektral}

Gunakan hasil kali dalam standar $\langle x,y\rangle=x^*y$ pada
$\C^n$, dengan $A^*=\overline A^{\,T}$ dan $\|x\|^2=x^*x$.
Matriks $U$ disebut \emph{uniter} jika $U^*U=UU^*=I$;
kolom-kolomnya merupakan basis ortonormal.
Versi realnya disebut matriks ortogonal, dengan $U^TU=I$.

\begin{teorema}[][Dekomposisi Schur][eigen:schur]
    Untuk setiap $A\in\Mat_n(\C)$, terdapat matriks uniter $U$
    sehingga $U^*AU=T$ segitiga atas. Entri diagonal $T$ adalah
    nilai eigen $A$, dihitung dengan multiplisitas aljabar.
\end{teorema}
\begin{proof}
    Pilih satu vektor eigen satuan $u_1$, lalu perluas menjadi basis
    ortonormal, misalnya dengan Gram--Schmidt. Jika $U_1$ mempunyai
    basis tersebut sebagai kolom-kolomnya, maka
    \[
        U_1^*AU_1=\begin{pmatrix}\lambda_1&w\\0&B\end{pmatrix}.
    \]
    Dengan induksi pada ukuran, terdapat uniter $V$ sehingga
    $V^*BV$ segitiga atas. Ambil $U=U_1\diag(1,V)$.
    Determinan matriks segitiga menunjukkan bahwa entri diagonalnya
    adalah akar polinom karakteristik, termasuk multiplisitasnya.
\end{proof}

Schur berlaku untuk semua matriks kompleks, termasuk yang tidak dapat
didiagonalkan. Dari segitiga atas, untuk setiap polinom $f$ diperoleh
\[
    \tr f(A)=\sum_{i=1}^n f(\lambda_i),\qquad
    \det f(A)=\prod_{i=1}^n f(\lambda_i).
\]
Secara khusus, $\tr A^k=\sum_i\lambda_i^k$ untuk $k\ge1$,
tanpa perlu mengasumsikan diagonalisasi.

\begin{definisi}[][Matriks normal dan Hermitian][eigen:normal]
    Matriks $A$ disebut \emph{normal} jika $A^*A=AA^*$,
    dan \emph{Hermitian} jika $A^*=A$.
    Matriks Hermitian disebut semidefinit positif jika $x^*Ax\ge0$
    untuk semua $x$, dan definit positif jika $x^*Ax>0$ untuk semua $x\ne0$.
\end{definisi}

\begin{teorema}[][Dekomposisi spektral][eigen:teorema-spektral]
    Matriks $A\in\Mat_n(\C)$ normal jika dan hanya jika
    \[
        A=U\diag(\lambda_1,\ldots,\lambda_n)U^*
    \]
    untuk suatu matriks uniter $U$.
    Khususnya, matriks Hermitian mempunyai nilai eigen real dan
    basis ortonormal yang terdiri atas vektor eigen.
\end{teorema}
\begin{proof}
    Jika $A=UDU^*$ dengan $D$ diagonal, maka
    $A^*A=UD^*DU^*=UDD^*U^*=AA^*$.

    Sebaliknya, gunakan Schur: $T=U^*AU$ segitiga atas dan normal.
    Kesamaan entri $(1,1)$ pada $TT^*=T^*T$ memberi
    \[
        \sum_{j=1}^n|t_{1j}|^2=|t_{11}|^2.
    \]
    Maka semua entri baris pertama di luar diagonal nol.
    Blok sisanya tetap normal dan segitiga atas; induksi menunjukkan
    bahwa $T$ diagonal.
    Jika $A=A^*$, untuk vektor eigen $v\ne0$ berlaku
    $\lambda=v^*Av/(v^*v)\in\R$.
\end{proof}

Untuk matriks simetris real, vektor eigen dapat dipilih real, lalu
diortonormalkan di setiap ruang eigen. Jadi terdapat matriks
ortogonal real $Q$ dengan $A=QDQ^T$.
Kondisi simetris kompleks $A^T=A$ berbeda dari Hermitian;
contoh pada latihan menunjukkan bahwa matriks simetris kompleks
bahkan dapat gagal didiagonalkan.

\subsection{Merekonstruksi matriks dari ruang eigennya}

Jika kolom-kolom $U$ adalah $u_1,\ldots,u_n$, bentuk spektral memberi
\[
    A=\sum_{i=1}^n\lambda_i u_iu_i^*.
\]
Dengan mengelompokkan nilai eigen yang sama, dapat ditulis
$A=\sum_\lambda\lambda P_\lambda$, dengan $P_\lambda$ proyektor
ortogonal ke $E_\lambda(A)$. Proyektor ini memenuhi
$P_\lambda^*=P_\lambda$, $P_\lambda^2=P_\lambda$, dan
$P_\lambda P_\mu=0$ jika $\lambda\ne\mu$.

\begin{exam}[Matriks simetris dari dua vektor eigen]
    Konstruksikan matriks simetris real $3\times3$ yang mempunyai
    nilai eigen $1,1,-1$, dengan $(1,1,1)^T$ dan $(2,2,1)^T$
    sebagai vektor eigen untuk $1$.
\end{exam}
\begin{solution}
    Kedua vektor bebas dan merentang bidang $x_1=x_2$.
    Komplemen ortogonalnya direntang oleh
    $u=(1,-1,0)^T/\sqrt2$, yang harus menjadi vektor eigen untuk $-1$.
    Matriks yang dicari adalah
    \[
        A=1(I-uu^T)-uu^T=I-2uu^T
        =\begin{pmatrix}0&1&0\\1&0&0\\0&0&1\end{pmatrix}.
    \]
    Ruang eigennya menentukan aksi pada seluruh ruang, sehingga
    matriks yang memenuhi syarat ini unik.
\end{solution}

\begin{prop}[Tanda bentuk kuadratik dan nilai eigen]
    \label{eigen:kuadratik}
    Jika $A$ Hermitian, maka $A$ semidefinit positif tepat ketika
    semua nilai eigennya nonnegatif, dan definit positif tepat ketika
    semua nilai eigennya positif. Selain itu,
    \[
        \lambda_{\min}\|x\|^2\le x^*Ax\le\lambda_{\max}\|x\|^2.
    \]
\end{prop}
\begin{proof}
    Tuliskan $A=UDU^*$ dan $y=U^*x$.
    Maka $x^*Ax=\sum_i\lambda_i|y_i|^2$ dan
    $\|y\|^2=\|x\|^2$. Ketaksamaan mengikuti dari nilai eigen terkecil
    dan terbesar. Syarat perlu untuk kepositifan diperoleh dengan
    mengambil $x$ berupa masing-masing vektor eigen satuan.
\end{proof}

Pada matriks normal, rank tepat sama dengan banyaknya nilai eigen
tak nol, termasuk multiplisitasnya, karena bentuk diagonal diperoleh
melalui perkalian invertibel. Pada matriks Hermitian, informasi
spektral juga menentukan tanda bentuk kuadratik dan memungkinkan
pemilihan basis ortonormal.

\section{Memilih Dekomposisi}
\label{eigen:peta-dekomposisi}

Keempat bentuk berikut menjawab kebutuhan yang berbeda.
Pada baris Jordan dan Schur, medan yang dipakai adalah $\C$.

\begin{center}
    \small
    \renewcommand{\arraystretch}{1.35}
    \begin{tabular}{>{\raggedright\arraybackslash}p{.18\linewidth}
        >{\raggedright\arraybackslash}p{.29\linewidth}
        >{\raggedright\arraybackslash}p{.42\linewidth}}
        \textbf{Bentuk} & \textbf{Syarat dan rumus} & \textbf{Informasi utama}\\\hline
        Diagonalisasi & Basis vektor eigen; $A=PDP^{-1}$ &
            Pangkat dan polinom matriks dihitung pada entri diagonal.\\
        Jordan & Semua matriks persegi kompleks; $A=PJP^{-1}$ &
            Ukuran rantai, polinom minimal, dan nulitas pangkat.\\
        Schur & Semua matriks persegi kompleks; $A=UTU^*$ &
            Bentuk segitiga dengan perubahan basis ortonormal.\\
        Spektral & Matriks normal; $A=UDU^*$ &
            Basis eigen ortonormal dan proyektor ortogonal.\\\hline
    \end{tabular}
\end{center}

\begin{enumerate}[label=\arabic*.,leftmargin=*,itemsep=3pt]
    \item Jika diberikan persamaan polinom untuk $A$, periksa apakah
        polinom itu mempunyai akar berulang sebelum mencari vektor eigen.
    \item Jika akar karakteristik berulang, hitung
        $\nul(A-\lambda I)$ dan bandingkan dengan multiplisitas aljabarnya.
    \item Jika diperlukan struktur yang lebih rinci, hitung nulitas
        pangkat $A-\lambda I$ untuk menentukan blok Jordan.
    \item Jika matriks Hermitian atau simetris real, manfaatkan
        ortogonalitas ruang eigen dan bentuk $\sum_i\lambda_i u_iu_i^*$.
    \item Untuk identitas trace yang melibatkan pangkat, Schur memberi
        akses ke nilai eigen tanpa mengasumsikan diagonalisasi.
    \item Jika dua matriks komutatif, periksa ruang eigen yang invarian;
        spektrum sederhana sering mengubah persoalan menjadi interpolasi.
\end{enumerate}

\section{Latihan Pilihan}
\label{dekomp:latihan}

Latihan berikut menggunakan diagonalisasi, bentuk Jordan, dan dekomposisi spektral.

\begin{latihan}\label{eigen:lat-komutatif}
    Misalkan $A$ invertibel dan mempunyai $n$ nilai eigen berbeda.
    Jika $AB=BA^{-1}$, buktikan bahwa $B^2$ dapat didiagonalkan.
    Apakah $B$ sendiri harus dapat didiagonalkan?
\end{latihan}

\begin{latihan}\label{eigen:lat-trace}
    Jika $A\in\Mat_n(\C)$ memenuhi $\tr A^k=0$ untuk
    $k=1,\ldots,n$, buktikan bahwa $A^n=0$.
    Buktikan pula arah sebaliknya.
\end{latihan}

\begin{latihan}\label{eigen:lat-real}
    Misalkan $A\in\Mat_2(\R)$ memenuhi $A^2=-I$.
    Buktikan bahwa $A$ serupa di atas $\R$ dengan
    $\begin{pmatrix}0&-1\\1&0\end{pmatrix}$.
    Jelaskan perbedaan diagonalisasinya di atas $\R$ dan $\C$.
\end{latihan}

\begin{latihan}\label{eigen:lat-rekonstruksi}
    Matriks simetris real $A$ mempunyai nilai eigen $1,2,3$.
    Diketahui $u=(-1,-1,1)^T$ dan $v=(1,-2,-1)^T$ merupakan
    vektor eigen untuk $1$ dan $2$. Tentukan vektor eigen untuk $3$
    dan hitung matriks $A$.
\end{latihan}

\begin{latihan}\label{eigen:lat-rank-diagonal}
    Buktikan bahwa $A\in\Mat_n(\C)$ dapat didiagonalkan jika dan
    hanya jika untuk setiap nilai eigen $\lambda$ berlaku
    \[
        \rank(A-\lambda I)=\rank(A-\lambda I)^2.
    \]
\end{latihan}

\begin{latihan}\label{eigen:lat-minimal-sama}
    Konstruksikan dua matriks $6\times6$ yang mempunyai polinom
    karakteristik dan polinom minimal sama, tetapi tidak serupa.
    Jelaskan mengapa setiap polinom yang menganulir salah satunya
    juga menganulir yang lain.
\end{latihan}

\begin{latihan}\label{eigen:lat-simetris-kompleks}
    Matriks $A=\begin{pmatrix}1&i\\i&-1\end{pmatrix}$ memenuhi $A^T=A$.
    Tentukan polinom karakteristik, polinom minimal, dan dimensi
    ruang eigennya. Apakah $A$ dapat didiagonalkan di atas $\C$?
\end{latihan}

\begin{latihan}\label{eigen:lat-frobenius}
    Misalkan $\lambda_1,\ldots,\lambda_n$ adalah nilai eigen
    $A=(a_{ij})\in\Mat_n(\C)$, dihitung dengan multiplisitas.
    Buktikan bahwa
    \[
        \sum_{i=1}^n|\lambda_i|^2\le\sum_{i,j=1}^n|a_{ij}|^2,
    \]
    dengan kesamaan tepat ketika $A$ normal.
\end{latihan}

\begin{latihan}\label{eigen:lat-komutator}
    Misalkan $A\in\Mat_n(\C)$ dapat didiagonalkan dan
    $\mathcal T:\Mat_n(\C)\to\Mat_n(\C)$ didefinisikan oleh
    $\mathcal T(X)=AX-XA$.
    Tentukan nilai eigen $\mathcal T$ dan buktikan bahwa pemetaan ini
    dapat didiagonalkan. Jika multiplisitas nilai eigen berbeda dari
    $A$ adalah $m_1,\ldots,m_s$, tentukan rank dan nulitas $\mathcal T$.
\end{latihan}

\begin{latihan}\label{eigen:lat-siklik}
    Vektor $v$ disebut siklik untuk $A\in\Mat_n(\C)$ jika
    $v,Av,\ldots,A^{n-1}v$ merupakan basis $\C^n$.
    \begin{enumerate}[label=(\alph*),nosep]
        \item Jika ada vektor siklik, buktikan $m_A=\chi_A$ dan bahwa
            setiap ruang eigen $A$ berdimensi satu.
        \item Jika $A$ mempunyai $n$ nilai eigen berbeda, konstruksikan
            satu vektor siklik.
    \end{enumerate}
\end{latihan}

\begin{latihan}\label{eigen:lat-blok}
    Untuk $A,B\in\Mat_n(\C)$, definisikan
    $M=\begin{pmatrix}A&B\\B&A\end{pmatrix}$.
    Nyatakan $\chi_M$ melalui $\chi_{A+B}$ dan $\chi_{A-B}$.
    Buktikan bahwa $M$ dapat didiagonalkan jika dan hanya jika
    $A+B$ dan $A-B$ masing-masing dapat didiagonalkan.
    Apakah diperlukan asumsi $AB=BA$?
\end{latihan}

\ifpembahasan
\section{Pembahasan Latihan}
\label{dekomp:pembahasan}

\subsection*{Latihan \ref{eigen:lat-komutatif}: kuadrat yang komutatif}
Dari $AB=BA^{-1}$ diperoleh $A^{-1}B=BA$.
Maka
\[
    AB^2=BA^{-1}B=B^2A.
\]
Karena $A$ mempunyai spektrum sederhana, \autoref{eigen:centralizer}
menjamin $B^2$ diagonal dalam basis vektor eigen $A$.
Matriks $B$ sendiri belum tentu dapat didiagonalkan. Contohnya,
\[
    A=\begin{pmatrix}2&0\\0&1/2\end{pmatrix},\qquad
    B=\begin{pmatrix}0&1\\0&0\end{pmatrix}.
\]
Relasi $AB=BA^{-1}$ berlaku, tetapi $B$ nilpoten tak nol.

\subsection*{Latihan \ref{eigen:lat-trace}: memisahkan nilai eigen tak nol}
Andaikan terdapat nilai eigen tak nol yang berbeda
$\mu_1,\ldots,\mu_s$, dengan multiplisitas aljabar $r_1,\ldots,r_s$.
Dari Schur dan hipotesis, untuk $k=1,\ldots,s$ berlaku
\[
    \sum_{j=1}^s r_j\mu_j^k=0.
\]
Matriks koefisiennya adalah matriks Vandermonde dengan setiap kolom
ke-$j$ dikalikan $\mu_j$. Determinannya tak nol karena semua $\mu_j$
berbeda dan tak nol. Akibatnya semua $r_j=0$, bertentangan dengan
multiplisitas yang positif.
Jadi seluruh nilai eigen $A$ nol, $\chi_A(t)=t^n$, dan $A^n=0$
oleh Cayley--Hamilton.
Sebaliknya, $A^n=0$ memaksa setiap nilai eigen nol; Schur memberi
$\tr A^k=0$ untuk semua $k\ge1$.

\subsection*{Latihan \ref{eigen:lat-real}: basis dari satu vektor dan bayangannya}
Pilih $v\ne0$ real. Vektor $v$ dan $Av$ bebas: jika $Av=cv$
dengan $c\in\R$, maka $-v=A^2v=c^2v$, yang mustahil.
Dalam basis $(v,Av)$, aksi $A$ adalah
\[
    v\longmapsto Av,\qquad Av\longmapsto-v.
\]
Jadi dengan $P=[v\ Av]$ diperoleh
$P^{-1}AP=\begin{pmatrix}0&-1\\1&0\end{pmatrix}$.
Di atas $\R$ tidak ada nilai eigen, sehingga diagonalisasi mustahil.
Di atas $\C$, kedua nilai eigen $i,-i$ berbeda, sehingga diagonalisasi ada.

\subsection*{Latihan \ref{eigen:lat-rekonstruksi}: proyektor ortogonal berbobot}
Vektor $u$ dan $v$ ortogonal. Vektor $w=(1,0,1)^T$ ortogonal terhadap
keduanya, sehingga menjadi vektor eigen untuk $3$.
Karena $\|u\|^2=3$, $\|v\|^2=6$, dan $\|w\|^2=2$,
\[
    A=\frac{uu^T}{3}+2\frac{vv^T}{6}+3\frac{ww^T}{2}
    =\frac16\begin{pmatrix}13&-2&5\\-2&10&2\\5&2&13\end{pmatrix}.
\]
Masing-masing suku mengatur aksi pada satu arah dalam basis ortonormal.

\subsection*{Latihan \ref{eigen:lat-rank-diagonal}: mendeteksi blok berukuran dua atau lebih}
Rank yang sama ekuivalen dengan
$\nul(A-\lambda I)=\nul(A-\lambda I)^2$.
Menurut \autoref{eigen:jordan-nulitas}, selisih kedua nulitas adalah
banyaknya blok Jordan untuk $\lambda$ yang berukuran setidaknya dua.
Jadi syarat tersebut untuk setiap nilai eigen berarti semua blok
berukuran satu, tepat syarat diagonalisasi.

\subsection*{Latihan \ref{eigen:lat-minimal-sama}: ukuran blok masih dapat berbeda}
Ambil
\[
    A=\diag\bigl(J_3(0),J_3(0)\bigr),\qquad
    B=\diag\bigl(J_3(0),J_2(0),J_1(0)\bigr).
\]
Keduanya mempunyai $\chi(t)=t^6$ dan $m(t)=t^3$.
Namun, $\nul A=2$ dan $\nul B=3$, sehingga keduanya tidak serupa.
Untuk kedua matriks, $f$ menganulirnya tepat ketika $t^3$ membagi $f$.
Jadi kesamaan seluruh polinom penganulir belum menentukan ukuran
semua blok Jordan.

\subsection*{Latihan \ref{eigen:lat-simetris-kompleks}: simetris tidak selalu Hermitian}
Perkalian langsung memberi $A^2=0$, sedangkan $A\ne0$.
Karena trace dan determinannya nol, $\chi_A(t)=t^2$.
Polinom minimalnya juga $t^2$, sebab $A$ tidak dianulir oleh $t$.
Rank $A$ adalah satu dan $\ker A=\operatorname{span}\{(-i,1)^T\}$.
Jadi satu-satunya nilai eigen adalah $0$, dengan multiplisitas
aljabar dua dan geometri satu; matriks tidak dapat didiagonalkan.
Tidak ada pertentangan dengan teorema spektral karena $A^*\ne A$.

\subsection*{Latihan \ref{eigen:lat-frobenius}: entri di atas diagonal Schur}
Untuk matriks persegi berukuran sama,
$\tr(XY)=\sum_{i,j}x_{ij}y_{ji}=\tr(YX)$.
Dengan dekomposisi Schur $A=UTU^*$,
\[
    \sum_{i,j}|a_{ij}|^2=\tr(A^*A)=\tr(T^*T)
    =\sum_i|\lambda_i|^2+\sum_{i<j}|t_{ij}|^2.
\]
Ini memberi ketaksamaan. Kesamaan berlaku tepat ketika semua entri
di atas diagonal $T$ nol, yaitu ketika $A$ dapat didiagonalkan
secara uniter. Oleh teorema spektral, ini ekuivalen dengan normalitas $A$.

\subsection*{Latihan \ref{eigen:lat-komutator}: basis matriks satuan}
Tuliskan $A=PDP^{-1}$, dengan $D=\diag(\lambda_1,\ldots,\lambda_n)$.
Untuk matriks satuan $E_{ij}$, yang hanya mempunyai entri $1$ pada
posisi $(i,j),$ berlaku
\[
    DE_{ij}-E_{ij}D=(\lambda_i-\lambda_j)E_{ij}.
\]
Karena $PE_{ij}P^{-1}$ untuk semua pasangan $(i,j)$ membentuk basis
$\Mat_n(\C)$, matriks-matriks ini merupakan basis vektor eigen
$\mathcal T$, dengan nilai eigen $\lambda_i-\lambda_j$.
Nilai eigen nol muncul tepat ketika kedua indeks berasal dari
nilai eigen $A$ yang sama. Dengan demikian
\[
    \nul\mathcal T=\sum_{j=1}^s m_j^2,\qquad
    \rank\mathcal T=n^2-\sum_{j=1}^s m_j^2.
\]

\subsection*{Latihan \ref{eigen:lat-siklik}: vektor siklik dan Vandermonde}
Jika $v,Av,\ldots,A^{n-1}v$ bebas, tidak ada polinom tak nol
berderajat kurang dari $n$ yang menganulir $A$: menerapkannya pada $v$
akan menghasilkan relasi linear tak trivial.
Maka $\deg m_A\ge n$; karena $m_A\mid\chi_A$ dan $\deg\chi_A=n$,
diperoleh $m_A=\chi_A$.
Dalam bentuk Jordan, pangkat setiap faktor dalam $m_A$ adalah ukuran
blok terbesar, sedangkan pangkatnya dalam $\chi_A$ adalah jumlah
ukuran semua blok untuk nilai eigen itu. Kesamaan keduanya memaksa
tepat satu blok per nilai eigen, sehingga setiap ruang eigen berdimensi satu.

Sebaliknya, jika nilai eigen $\lambda_1,\ldots,\lambda_n$ berbeda,
pilih basis vektor eigen $v_1,\ldots,v_n$ dan ambil $v=\sum_i v_i$.
Koordinat $v,Av,\ldots,A^{n-1}v$ terhadap basis ini membentuk matriks
Vandermonde dengan baris ke-$i$ sama dengan
$(1,\lambda_i,\ldots,\lambda_i^{n-1})$.
Determinan tak nol, sehingga $v$ siklik.

\subsection*{Latihan \ref{eigen:lat-blok}: memisahkan jumlah dan selisih}
Matriks
\[
    U=\frac1{\sqrt2}\begin{pmatrix}I&I\\I&-I\end{pmatrix}
\]
uniter dan memenuhi
\[
    U^*MU=\begin{pmatrix}A+B&0\\0&A-B\end{pmatrix}.
\]
Jadi $\chi_M(t)=\chi_{A+B}(t)\chi_{A-B}(t)$.
Jika kedua blok dapat didiagonalkan, gabungan perubahan basisnya
mendiagonalkan matriks diagonal blok, dan karenanya juga $M$.
Sebaliknya, jika $M$ dapat didiagonalkan, suatu polinom dengan faktor
linear berbeda menganulir $M$ dan kedua bloknya. Maka kedua blok
dapat didiagonalkan oleh \autoref{eigen:squarefree}.
Tidak diperlukan $AB=BA$; identitas blok di atas berlaku tanpa asumsi itu.

\fi
